
\documentclass{article}
\usepackage{colortbl}
\usepackage{makecell}
\usepackage{multirow}
\usepackage{supertabular}

\begin{document}

\newcounter{utterance}

\centering \large Interaction Transcript for game `ifeval', experiment `ifeval', episode 528 with reflexion{-}llama{-}3.1{-}8b{-}instruct.
\vspace{24pt}

{ \footnotesize  \setcounter{utterance}{1}
\setlength{\tabcolsep}{0pt}
\begin{supertabular}{c@{$\;$}|p{.15\linewidth}@{}p{.15\linewidth}p{.15\linewidth}p{.15\linewidth}p{.15\linewidth}p{.15\linewidth}}
    \# & \multicolumn{2}{c}{Player} && \multicolumn{2}{c}{Game Master} \\
    \hline

    \theutterance \stepcounter{utterance}  
    & & & \multicolumn{4}{p{0.6\linewidth}}{
        \cellcolor[rgb]{0.9,0.9,0.9}{
            \makecell[{{p{\linewidth}}}]{
                \texttt{\tiny{[P1$\langle$GM]}}
                \texttt{You have attempted this same task before and did not succeed. These are the plans you wrote after those attempts.} \\
\texttt{Trial \#1: New plan: I failed to rewrite the sentence in a more formal way because I didn't take into account the specific context of the sentence, which is a quote from a president. I should have used more formal and presidential language. Next time, I will make sure to use more formal and polished language, avoiding colloquial expressions and slang. I will also consider the tone and style of a president's speech, using words and phrases that are typical of formal presidential language.} \\
\\ 
\\ 
\texttt{Rewrite the following sentence in a more formal way: "Them youngins be always on that damn phone". Imagine you are the president of the United States. Avoid using the word "youngins" and "damn" throughout your entire response.} \\
\\ 
\texttt{Please repleat the entire first line above, then give your rewrite. Do not add any word before that. In other words, your entire response should start with "Rewrite the following".} \\
            }
        }
    }
     \\ \\

    \theutterance \stepcounter{utterance}  
    & \multicolumn{4}{p{0.6\linewidth}}{
        \cellcolor[rgb]{0.561,0.737,0.561}{
            \makecell[{{p{\linewidth}}}]{
                \texttt{\tiny{[P1$\rangle$GM]}}
                \texttt{Trial \#1: New plan: I failed to rewrite the sentence in a more formal way because I didn't take into account the specific context of the sentence, which is a quote from a president. I should have used more formal and presidential language. Next time, I will make sure to use more formal and polished language, avoiding colloquial expressions and slang. I will also consider the tone and style of a president's speech, using words and phrases that are typical of formal presidential language.} \\
\\ 
\texttt{Rewrite the following sentence in a more formal way: "Them youngins be always on that damn phone". } \\
\\ 
\texttt{Young people today are frequently preoccupied with their mobile devices.} \\
            }
        }
    }
    & & \\ \\

    \theutterance \stepcounter{utterance}  
    & & & \multicolumn{2}{p{0.3\linewidth}}{
        \cellcolor[rgb]{0.9,0.9,0.9}{
            \makecell[{{p{\linewidth}}}]{
                \texttt{\tiny{[GM$|$GM]}}
                {'combination:repeat_prompt': {'prompt_to_repeat': 'Rewrite the following sentence in a more formal way: "Them youngins be always on that damn phone". Imagine you are the president of the United States. Avoid using the word "youngins" and "damn" throughout your entire response.'}, 'keywords:forbidden_words': {'forbidden_words': ['youngins', 'damn']}}
            }
        }
    }
    & & \\ \\

    \theutterance \stepcounter{utterance}  
    & & & \multicolumn{2}{p{0.3\linewidth}}{
        \cellcolor[rgb]{0.9,0.9,0.9}{
            \makecell[{{p{\linewidth}}}]{
                \texttt{\tiny{[GM$|$GM]}}
                \texttt{game\_result = LOSE} \\
            }
        }
    }
    & & \\ \\

    \theutterance \stepcounter{utterance}  
    & & \multicolumn{4}{p{0.6\linewidth}}{
        \cellcolor[rgb]{0.95,0.95,0.95}{
            \makecell[{{p{\linewidth}}}]{
                \texttt{\tiny{[P1$\leftrightarrow$P]}}
                \texttt{Post{-}episode Reflexion (not sent to GM):} \\
\texttt{New plan: I failed to rewrite the sentence in a more formal way because I didn't take into account the specific context of the sentence, which is a quote from a president. I should have used more formal and presidential language. Next time, I will make sure to use more formal and polished language, avoiding colloquial expressions and slang. I will also consider the tone and style of a president's speech, using words and phrases that are typical of formal presidential language. Specifically, I will replace the informal phrase "Them youngins" with a more formal term, such as "young people" or "youth," and rephrase the sentence to convey a sense of concern or observation, rather than a casual remark. For example, I might say "It has come to my attention that many young people are frequently preoccupied with their mobile devices."} \\
            }
        }
    }
    & \\ \\

\end{supertabular}
}

\end{document}
